\exercisesection{MULTIPLICITIES}

\begin{enumerate}
\def\labelenumi{\arabic{enumi})}
\tightlist
\item
  Let \(\Gamma \subset \mathbf{N}\) be a submonoid of \(\mathbf{N}\) such that \(\mathbf{N} - \Gamma\) is a finite set. Let \((a_1, \ldots, a_r)\), with \(0 < a_1 < \cdots < a_r\), be a minimal system of generators of \(\Gamma\), and let \(k\) be a field. Set \(\mathbf{A} = k[\Gamma]\).
\end{enumerate}

\begin{enumerate}
\def\labelenumi{\alph{enumi})}
\item
  Show that A is isomorphic to the sub-k-algebra of \(k[\mathbf{T}]\) generated by \(\mathbf{T}^{a_1}\), \(\ldots\), and \(\mathbf{T}^{a_r}\).
\item
  Let m be the maximal ideal of A generated by \(T^{a_{1}}\), \(\ldots\), and \(T^{a_{r}}\); let \(A_{m}\) be the local ring of A at m, and set \(q = mA_{m}\). Show that we have \(\dim_{k}q/q^{2} = r\) and \(e_{q}(A_{m}) = a_{1}\).
\end{enumerate}

\begin{enumerate}
\def\labelenumi{\arabic{enumi})}
\setcounter{enumi}{1}
\tightlist
\item
  Let \(k\) be a field, \(m\) an integer \(\geqslant 1\), and I an ideal of the polynomial ring \(k[\mathbf{T}_{1},\ldots,\mathbf{T}_{m}]\) generated by monomials represented by points \(\mathbf{M}_{1},\ldots,\mathbf{M}_{s}\) of the “positive quadrant” \(\mathbf{R}_{+}^{m}\) of \(\mathbf{R}^{m}\). Denote by \(\mathcal{N}\) the convex hull in \(\mathbf{R}_{+}^{m}\) of the union of the sets \(\mathbf{M}_{i}+\mathbf{R}_{+}^{m}\) for \(1\leqslant i\leqslant s\).
\end{enumerate}

\begin{enumerate}
\def\labelenumi{\alph{enumi})}
\tightlist
\item
  Set \(m = (T_1, \ldots, T_m)\), \(A = k[T_1, \ldots, T_m]_m\), and \(q = I.A\). Show that \(A/q\) has finite length if and only if the volume of \(\mathbf{R}_+^m - \mathcal{N}\) is finite.
\item
  Prove the equality \(e_q(A) = m! \operatorname{Vol}(\mathbf{R}_+^m - \mathcal{N})\).
\end{enumerate}

\begin{enumerate}
\def\labelenumi{\arabic{enumi})}
\setcounter{enumi}{2}
\tightlist
\item
  Let \(k\) be a field, and let \(f_{1},...,f_{r}\) be homogeneous elements of the graded ring \(k[X_{1},...,X_{n}]\) forming a complete intersection sequence. One says that the quotient ring
\end{enumerate}

\[
\mathbf {H} = k [ \mathrm{X} _ {1},..., \mathrm{X} _ {n} ] / (f _ {1},..., f _ {r})
\]

is a graded complete intersection ring. Let I be a nonzero graded ideal of the ring \(\mathbf{H} = k[\mathbf{X}_1, \ldots, \mathbf{X}_n] / (f_1, \ldots, f_r)\) endowed with the quotient grading. Show that one has \(\dim(\mathbf{I}) = \dim(\mathbf{H})\) . Deduce that one has either \(\dim(\mathbf{H}/\mathbf{I}) < \dim(\mathbf{H})\) , or \(\dim(\mathbf{H}/\mathbf{I}) = \dim(\mathbf{H})\) and \(c_{\mathbf{M}/\mathbf{I}} < c_{\mathbf{M}}\) ( \(\S 4, \text{n}^{\circ} 2\) ). (One will consider the ideal b formed by the \(s \in \mathbf{H}\) such that \(s.\mathbf{I} = \{0\}\) and we will show that if one has \(\dim(\mathbf{I}) < \dim(\mathbf{H})\) , then b contains a non-zero-divisor element, by examining the position of b relative to the minimal prime ideals of H and to \(\mathbf{H}_{+} = \mathbf{H}_{\geqslant 1}\) .)

\begin{enumerate}
\def\labelenumi{\arabic{enumi})}
\setcounter{enumi}{3}
\tightlist
\item
  Let A be a Noetherian local ring, with maximal ideal m.
\end{enumerate}

\begin{enumerate}
\def\labelenumi{\alph{enumi})}
\item
  Suppose that \(\mathrm{gr}_{\mathfrak{m}}(\mathrm{A})\) is a graded complete intersection ring (cf. Exercise 3). Show that the completion \(\widehat{A}\) of A is a complete intersection, that is, the quotient of a regular local ring by an ideal generated by a completely secant sequence.
\item
  Let \(x \in \mathfrak{m}\), and let \(\xi\) be the image of \(x\) in \(\mathfrak{m}^{\nu}/\mathfrak{m}^{\nu+1}\), where \(\nu\) is such that \(x \in \mathfrak{m}^{\nu}\) and \(x \notin \mathfrak{m}^{\nu+1}\). Show that we have \(e_{\mathfrak{m}}(\mathbf{A}/x\mathbf{A}) = \nu e_{\mathfrak{m}}(\mathbf{A})\) if and only if \(\xi\) is not a zero divisor in \(\mathrm{gr}_{\mathfrak{m}}(\mathbf{A})\), and that in this case one has \(\mathrm{gr}_{\mathfrak{m}}(\mathbf{A}/x\mathbf{A}) = \mathrm{gr}_{\mathfrak{m}}(\mathbf{A})/\xi\mathrm{gr}_{\mathfrak{m}}(\mathbf{A})\) and \(x\) is not a zero divisor in A. (Let I be the homogeneous ideal of \(\mathrm{gr}_{\mathfrak{m}}(\mathbf{A})\) consisting of the elements \(\eta \in \mathrm{gr}_{\mathfrak{m}}(\mathbf{A})\) such that \(\xi\eta = 0\). Consider the ideal \(\mathrm{N}_n = \mathfrak{m}^{\nu+n}:x\) of A and the map from \(\mathrm{I}_n\) into \(\mathrm{N}_{n+1}/\mathfrak{m}^{n+1}\) which sends \(\eta \in \mathrm{I}_n\) to the class modulo \(\mathfrak{m}^{n+1}\) of an element \(y \in A\) lifting \(\eta\). We will show that it is injective, and from this we will deduce the inequality
\end{enumerate}

\[
\operatorname{long} _ {\mathbf {A} / \mathfrak {m}} (\mathbf {I} _ {n}) \leqslant \operatorname{long} _ {\mathbf {A} / \mathfrak {m}} (\mathbf {N} _ {n + 1} / \mathfrak {m} ^ {n + 1});\tag{i}
\]

Moreover, we will prove the equality

\[
(1 - \mathrm{T} ^ {\mathrm{v}}) \mathrm{H} _ {\mathrm{A}} ^ {(1)} = \mathrm{T} ^ {- \mathrm{v}} \mathrm{H} _ {\mathrm{A} / x \mathrm{A}} ^ {(1)} - \sum_ {n \geqslant \mathrm{v}} \operatorname{long} _ {\mathrm{A} / \mathfrak {m}} (\mathrm{N} _ {n} / \mathfrak {m} ^ {n}). \mathrm{T} ^ {n}.\tag{ii}
\]

Finally, we note that \(\xi\) is a zero divisor if and only if I \(\neq\) \{0\}, and using the previous exercise, we deduce from (i) that \(\sum_{n\geqslant 0}\mathrm{long}_{\mathrm{A / m}}(\mathbf{N}_{n + 1} / \mathfrak{m}^{n + 1})\) \(\mathbf{T}^n = \frac{\mathbf{S}(\mathbf{T})}{(1 - \mathbf{T})^\nu}\), with \(\mathbf{S}(1)\neq 0\), and from (ii) the inequality \(e_{\mathfrak{m}}(\mathbf{A} / x\mathbf{A}) > e_{\mathfrak{m}}(\mathbf{A})\). c) Let \((x_{1},\dots,x_{s})\) be a sequence of elements of the maximal ideal m. Suppose we have \(x_{i}\in \mathfrak{m}^{\nu_i}\) and \(x_{i}\notin \mathfrak{m}^{\nu_i + 1}\), and let \(\xi_{i}\) be the class of \(x_{i}\) in \(\mathfrak{m}^{\nu_i} / \mathfrak{m}^{\nu_i + 1}\). Show that we have \(e_{\mathfrak{m}}(\mathbf{A} / x)\geqslant e_{\mathfrak{m}}(\mathbf{A})\nu_{1}\dots \nu_{s}\), and that equality holds if and only if the sequence \((\xi_1,\dots,\xi_s)\) is completely secant in the ring \(\mathrm{gr}_{\mathfrak{m}}(\mathbf{A})\). In this case we have an isomorphism of graded rings from \(\mathrm{gr}_{\mathfrak{m}}(\mathbf{A} / x)\) with \((\mathrm{gr}_{\mathfrak{m}}(\mathbf{A}))/\xi\), where \(\mathbf{x} = \sum_{i}x_{i}\mathbf{A}\) and \(\xi = \sum_{i}\xi_{i}\mathrm{gr}_{\mathfrak{m}}(\mathbf{A})\).

\begin{enumerate}
\def\labelenumi{\arabic{enumi})}
\setcounter{enumi}{4}
\tightlist
\item
  Let \(k\) be a field. Set \(\mathbf{R} = k[[\mathrm{X},\mathrm{Y},\mathrm{Z}]]\) and consider the local ring \(\mathbf{A} = \mathbf{R}/(\mathbf{XY},\mathbf{X}\mathbf{Z})\) with maximal ideal m. Show that one has
\end{enumerate}

\[
\mathrm{H} _ {\mathrm{A}} = (1 - \mathrm{T}) (1 - \mathrm{T} - \mathrm{T} ^ {2}) \mathrm{H} _ {\mathrm{R}} = 1 + \sum_ {n \geqslant 1} (n + 2) \mathrm{T} ^ {n}.
\]

Hence we have \(\dim(A) = 2\), \(\dim_k(m/m^2) = 3\), and \(e_m(A) = 1\). (We have \(H_A = P_G\), where \(G = gr_m(A) = k[X, Y, Z]/(XY, XZ)\). We shall prove the existence of an exact sequence of graded H-modules

\[
0 \rightarrow \mathrm{H} (- 3) \rightarrow \mathrm{H} (- 2) \oplus \mathrm{H} (- 2) \rightarrow \mathrm{H} \rightarrow \mathrm{G} \rightarrow 0
\]

\(\text{where } \mathrm{H} = k[\mathrm{X},\mathrm{Y},\mathrm{Z}].\)

\begin{enumerate}
\def\labelenumi{\arabic{enumi})}
\setcounter{enumi}{5}
\item
  Let \(k\) be a field, A a local, complete, reduced \(k\)-algebra of dimension 1, and \(m\) its maximal ideal. Show that the minimum number of generators of an ideal of A is at most \(e_m(A)\). There exists \(X \in A\) such that the injection of \(k[[X]]\) into A makes A into a \(k[[X]]\)-free module of rank \(e_m(A)\).
\item
  Let A be a local ring with maximal ideal m, and let M be a finitely generated A-module, \((x_{1}, \ldots, x_{d})\) a sequence of elements of m that is secant for the A-module M, and x the ideal \(x_{1}A + \cdots + x_{d}A\).
\end{enumerate}

\begin{enumerate}
\def\labelenumi{\alph{enumi})}
\tightlist
\item
  Suppose the A-module M is faithful. Then the quotient of \(\mathrm{gr}_{x}(\mathbf{A})\) by its nilradical n is isomorphic to \((\mathbf{A}/\mathfrak{m})[\mathbf{X}_{1}, \ldots, \mathbf{X}_{d}]\), and \((\mathrm{gr}_{x}(\mathbf{M}))_{n}\) is a \((\mathrm{gr}_{x}(\mathbf{A}))_{n}\)-module of finite length \(e_{x}(\mathbf{M})\).
\item
  In the general case, \(\mathrm{gr}_{x}(\mathbf{A})\) has a unique minimal prime ideal n, and \((\mathrm{gr}_{x}(\mathbf{M}))_{n}\) is a \(\mathrm{gr}_{x}(\mathbf{A})\)-module of finite length \(e_{x}(\mathbf{M})\).
\end{enumerate}

\begin{enumerate}
\def\labelenumi{\arabic{enumi})}
\setcounter{enumi}{7}
\tightlist
\item
  Let A be a ring and a an ideal of A. Set \(a^{i} = A\) for \(i \leqslant 0\) and consider the A-algebra \(\mathcal{A}(\mathfrak{a}) = \sum_{i \in \mathbf{Z}} \mathfrak{a}^{-i} v^{i} \subset \mathrm{A}[v, v^{-1}]\) and the A-algebra \(\mathrm{P}(\mathfrak{a}) = \sum_{i \geqslant 0} \mathfrak{a}^{i} v^{i} \subset \mathrm{A}[v]\). These are two graded A-algebras having the same total ring of fractions. Show that for an element \(a \in A\) the following conditions are equivalent:
\end{enumerate}

\begin{enumerate}
\def\labelenumi{(\roman{enumi})}
\item
  The element \(av\) of the \(\mathcal{A}(\mathfrak{a})\)-algebra \(\mathbf{A}[v, v^{-1}]\) is integral over \(\mathcal{A}(\mathfrak{a})\).
\item
  The element av of the P(a)-algebra A{[}v{]} is integral over P(a).
\item
  There exists an integral dependence relation:
\end{enumerate}

\[
a ^ {k} + b _ {1} a ^ {k - 1} + \dots + b _ {k} = 0 \quad \text { with } \quad b _ {j} \in \mathfrak {a} ^ {j}.
\]

One says that \(a\) is integral over the ideal \(\mathfrak{a}\).

\begin{enumerate}
\def\labelenumi{\arabic{enumi})}
\setcounter{enumi}{8}
\tightlist
\item
  \begin{enumerate}
  \def\labelenumii{\alph{enumii})}
  \tightlist
  \item
    Let A and a be as above, and let b be an ideal of A containing a. Show that the following conditions are equivalent:
  \end{enumerate}
\end{enumerate}

\begin{enumerate}
\def\labelenumi{(\roman{enumi})}
\item
  Every element of b is integral over a.
\item
  The \(\mathcal{A}(\mathfrak{a})\)-graded algebra \(\mathcal{A}(\mathfrak{b})\) is integral.
\item
  The graded \(P(\mathfrak{a})\)-algebra \(P(\mathfrak{b})\) is integral.
\item
  There exists an integer \(k \geqslant 0\) such that \(\mathfrak{a}\mathfrak{b}^{k} = \mathfrak{b}^{k+1}\) and therefore \(\mathfrak{a}^{n}\mathfrak{b}^{k} = \mathfrak{b}^{k+n}\) for all positive integers \(n\). We say that \(\mathfrak{b}\) is integral over \(\mathfrak{a}\).
\end{enumerate}

\begin{enumerate}
\def\labelenumi{\alph{enumi})}
\setcounter{enumi}{1}
\item
  Show that there exists a largest ideal among those that are integral over a. It is denoted \(\overline{a}\) and is called the integral closure of a in A. Verify that if \(a \subset b\), then \(\overline{a} \subset \overline{b}\) and that \(\overline{\overline{a}} = \overline{a}\). Show that if the integral closure \(\overline{P(a)}\) is a finitely generated \(P(a)\)-module, then \(\mathfrak{a}\,\overline{\mathfrak{a}^{k}} = \overline{\mathfrak{a}^{k+1}}\) for k sufficiently large.
\item
  Show that if \(a \in A\) satisfies the integral dependence relation
\end{enumerate}

\[
a ^ {k} + b _ {1} a ^ {k - 1} + \dots + b _ {k} = 0 \quad \text { with } \quad k \geqslant 1 \quad \text { and } \quad b _ {i} \in \mathfrak {a} ^ {i},
\]

then \(\mathfrak{a}(\mathfrak{a} + \mathbf{A}a)^{k-1} = (\mathfrak{a} + \mathbf{A}a)^k\). Deduce from this that the integral closure of \(\mathfrak{a}\) in A coincides with the set of elements of \(\mathbf{A}\) which are integral over \(\mathfrak{a}\).

10) Let A be a ring.

\begin{enumerate}
\def\labelenumi{\alph{enumi})}
\item
  Let B be an integral A-algebra. Show that for every element a of A and every ideal \(\mathfrak{a}\) of A, the element a is integral over \(\mathfrak{a}\) if and only if a is integral over the ideal \(\mathfrak{a}B\) generated by \(\mathfrak{a}\) in B.
\item
  Assuming A is Noetherian and reduced, denote \(\mathfrak{p}_{1},\ldots,\mathfrak{p}_{r}\) the minimal prime ideals of A and consider the natural injection:
\end{enumerate}

\[
\mathrm{A} \rightarrow \prod_ {1 \leqslant i \leqslant r} \mathrm{A} / \mathfrak {p} _ {i}.
\]

Show that it makes \(\prod_{1\leqslant i\leqslant r}\mathrm{A} / \mathfrak{p}_i\) into an integral A-algebra.

\begin{enumerate}
\def\labelenumi{\alph{enumi})}
\setcounter{enumi}{2}
\item
  Deduce from this that, given a Noetherian ring A, for an element a of A to be integral over an ideal \(\mathfrak{a}\), it is necessary and sufficient that, for every minimal prime ideal p of A, the image of a in A/p be integral over the ideal \((\mathfrak{a} + p)/p\).
\item
  Show that if \(\overline{a} = \overline{b}\), the ideals \(a\) and \(b\) have the same minimal associated prime ideals.
\end{enumerate}

\begin{enumerate}
\def\labelenumi{\arabic{enumi})}
\setcounter{enumi}{10}
\tightlist
\item
  Let A be a Noetherian ring, and let a, b be two ideals of A such that Supp(a) = Supp(b) = Spec(A). Show that the following conditions are equivalent:
\end{enumerate}

\begin{enumerate}
\def\labelenumi{(\roman{enumi})}
\tightlist
\item
  There exists a finitely generated A-module M such that \(\operatorname{Supp}(\mathbf{M}) = \operatorname{Spec}(\mathbf{A})\) and such that \(bM \subset aM\).\\
\item
  We have \(a, b \subset \overline{a}\).
\end{enumerate}

\begin{enumerate}
\def\labelenumi{\arabic{enumi})}
\setcounter{enumi}{11}
\tightlist
\item
  Let A be a Noetherian ring, and let a, b, and c be three ideals of A such that c is contained in the radical of A.
\end{enumerate}

\begin{enumerate}
\def\labelenumi{\alph{enumi})}
\item
  Show that if one has the inclusion \(\overline{b + c.a} \subset \overline{a}\), then we have \(\overline{b} \subset \overline{a}\).
\item
  Hence, if A is local, and if one has \(b \subset a\) and \(\overline{b} = \overline{a}\), there exists at least one ideal \(b_{1} \subset b\) such that
\end{enumerate}

\[
\overline {{{\mathbf {b} _ {1}}}} = \overline {{{\mathbf {a}}}},
\]

\(\beta)\) \(\mathbf{b}_{1}\) is minimal for this property, that is, if \(\mathbf{b}_{1}^{\prime}\subset \mathbf{b}_{1}\) and \(\overline{\mathbf{b}_1'} = \overline{\mathbf{b}_1}\), then \(\mathbf{b}_1' = \mathbf{b}_1\).

\begin{enumerate}
\def\labelenumi{\arabic{enumi})}
\setcounter{enumi}{12}
\item
  Let A be a local ring and let a and b be two ideals of A distinct from A. Show that if a ⊂ b, then one has \(\overline{a} = \overline{b}\) if and only if the graded algebra gr \(_b(A)\) is integral over the graded algebra gr \(_a(A)\) . Deduce that if a = (x \(_1\) , \ldots, x \(_d\) ) ⊂ b, where d = dim(A), we have \(\overline{a} = \overline{b}\) if and only if the initial forms \(\xi_{1}\) , \ldots, \(\xi_{d}\) in gr \(_b(A)\) of the x \(_i\) are of degree 1 and form a maximal secant sequence for gr \(_b(A)\) .
\item
  Let A be a local ring, m its maximal ideal, and q an ideal of A, distinct from A, such that A/q has finite length. Suppose that A contains a subfield k such that A = k + m.
\end{enumerate}

\begin{enumerate}
\def\labelenumi{\alph{enumi})}
\tightlist
\item
  Show that there exist elements \(x_{1}, \ldots, x_{d}\) of q, where \(d = \dim(A)\), with \(x_{i} \in q^{\delta_{i}}\) such that, for every sufficiently large integer v, one has \(q^{v} = x_{1} \cdot q^{v - \delta_{1}} + \cdots + x_{d} \cdot q^{v - \delta_{d}}\).
\item
  If one assumes that the field A/m is infinite, one can find an ideal a generated by d elements of q, where \(d = \dim(A)\), such that \(\overline{a} = \overline{q}\) (One may use § 7, no. 5.)
\end{enumerate}

\begin{enumerate}
\def\labelenumi{\arabic{enumi})}
\setcounter{enumi}{14}
\tightlist
\item
  Let A be a local, Noetherian, complete ring with maximal ideal m. Set \(d = \dim(A)\). Assume that A contains a subfield k such that \(A = k + m\).
\end{enumerate}

Show that one can find \(x_{1},\ldots,x_{d}\) in m such that:

the morphism \(\varphi:k[[X_{1},...,X_{d}]]\to A\) defined by \(\varphi(X_{i})=x_{i}\) makes A an integral \(k[[X_{1},...,X_{d}]]\)-algebra;

\(\beta)\) denoting \(\mathfrak{m}_0 = (X_1, \ldots, X_d)\) the maximal ideal of \(k[[X_1, \ldots, X_d]]\), one has \(\overline{\mathfrak{m}_0A} = \mathfrak{m}\) and therefore \(e_{\mathfrak{m}_0\mathbf{A}}(\mathbf{A}) = e_{\mathfrak{m}}(\mathbf{A})\).

Show that if \(x_{1},\ldots,x_{d}\) satisfy condition \(\alpha\), condition \(\beta\) is satisfied if and only if the initial forms \(\xi_{1},\ldots,\xi_{d}\) of \(x_{i}\) in \(\mathrm{gr}_{\mathfrak{m}}(\mathbf{A})\) form a maximal secant sequence.

\begin{enumerate}
\def\labelenumi{\arabic{enumi})}
\setcounter{enumi}{15}
\item
  Let A be a Noetherian local ring and \(x = (x_1, \ldots, x_d)\) an ideal generated by a maximal secant sequence for A. Show that we have the inequality \(e_x(A) \leqslant \text{long}(A/x)\) and that equality holds only if the homomorphism \((A/x)[T_1, \ldots, T_d] \xrightarrow{\varphi} \text{gr}_x(A)\) defined by \(\varphi(T_i) = \xi_i\), where \(\xi_i\) is the class of \(x_i\) modulo \(x^2\), is an isomorphism (which means that \((x_1, \ldots, x_d)\) is completely secant).
\item
  Let A be a Noetherian local ring possessing a completely secant sequence of length \(d = \dim(A)\), let \(\mathfrak{a}\) be the ideal generated by such a sequence, and let \(f \in A\) be an element integral over \(\mathfrak{a}\) but not belonging to it (for example \(A = k[[X_1, ..., X_d]]\), where k is a field, \(\mathfrak{a} = (X_1^k, ..., X_d^k)\), and \(f = X_1^{k-1}X_2\)). Let \(\mathfrak{a}' = \mathfrak{a} + fA\). Show that we have the inequality \(e_{\mathfrak{a}'}(A) > \text{long}(A/\mathfrak{a}')\).
\item
  Let A be a ring and \(\mathfrak{a}\) an ideal of A. For every element \(a \in \mathbf{A}\), denote by \(\nu_{\mathfrak{a}}(a)\) the upper bound of the set of integers \(\nu \in \mathbf{N}\) such that \(a \in \mathfrak{a}^{\nu}\).
\end{enumerate}

\begin{enumerate}
\def\labelenumi{\alph{enumi})}
\tightlist
\item
  Show that if \(a \notin \bigcap_{k=0}^{\infty} \mathfrak{a}^k\) the sequence \(\left(\frac{\nu_{\mathfrak{a}}(a^k)}{k}\right)\) is convergent. We denote \(\overline{\nu}_{\mathfrak{a}}(a)\) the limit. If
\end{enumerate}

\(a \in \bigcap_{k=0}^{\infty} \mathfrak{a}^{k}\) one sets \(\overline{\nu}_{\mathfrak{a}}(a) = +\infty\).

\begin{enumerate}
\def\labelenumi{\alph{enumi})}
\setcounter{enumi}{1}
\tightlist
\item
  Show that, for all \(a \in \mathbf{A}\), one has \(\overline{\nu}_{\mathfrak{a}}(a) = \overline{\nu}_{\mathfrak{a}}(a)\) and deduce from this that the following conditions are equivalent:
\end{enumerate}

\begin{enumerate}
\def\labelenumi{(\roman{enumi})}
\item
  \(\overline{\nu}_{\mathfrak{a}}(a)\geqslant 1,\)
\item
  \(a \in \overline{\mathfrak{a}}\).
\end{enumerate}

(Study the monotonicity of the sequence considered at a).)

\begin{enumerate}
\def\labelenumi{\arabic{enumi})}
\setcounter{enumi}{18}
\tightlist
\item
  Let A be a Noetherian ring and q an ideal of A such that A/q has finite length and q is contained in the radical of A. Show that we have
\end{enumerate}

\[
e _ {\mathfrak {q}} (\mathrm{A}) = e _ {\overline {{\mathfrak {q}}}} (\mathrm{A}).
\]

\begin{enumerate}
\def\labelenumi{\arabic{enumi})}
\setcounter{enumi}{19}
\tightlist
\item
  Let A be a Noetherian ring, M a finitely generated A-module, and \(q_{1}, \ldots, q_{k}\) ideals of A such that \(M/q_{i}M\) has finite length for every i. We assume \(M \neq 0\) and set \(d = \dim_{\mathbf{A}}(\mathbf{M})\). a) Show that \(M/q_{1}^{n_{1}} \ldots q_{k}^{n_{k}}M\) is of finite length for every k-tuple of positive integers \((n_{1}, \ldots, n_{k})\).
\end{enumerate}

\begin{enumerate}
\def\labelenumi{\alph{enumi})}
\setcounter{enumi}{1}
\tightlist
\item
  Consider the graded ring H of type \(\mathbf{N}^k\) such that
\end{enumerate}

\[
\mathrm{H} _ {n _ {1}, \dots , n _ {k}} = \mathfrak {q} _ {1} ^ {n _ {1}} \mathfrak {q} _ {2} ^ {n _ {2}} \dots \mathfrak {q} _ {k} ^ {n _ {k}} / \mathfrak {q} _ {1} ^ {n _ {1} + 1} \mathfrak {q} _ {2} ^ {n _ {2}} \dots \mathfrak {q} _ {k} ^ {n _ {k}},
\]

and the graded H-module N such that

\[
 \mathrm{N} _ {n _ {1}, \dots , n _ {k}} = \mathrm{q} _ {1} ^ {n _ {1}} \mathrm{q} _ {2} ^ {n _ {2}} \dots \mathrm{q} _ {k} ^ {n _ {k}} \mathrm{M} / \mathrm{q} _ {1} ^ {n _ {1} + 1} \mathrm{q} _ {2} ^ {n _ {2}} \dots \mathrm{q} _ {k} ^ {n _ {k}} \mathrm{M};
\]

show that N is a finitely generated graded H-module generated by \(N_{0,\ldots,0}\) and a finite number of elements whose degrees are vectors of the canonical basis of \(\mathbf{Z}^{k}\).

From this, deduce that there exists a polynomial \(\mathbf{Q}(\mathrm{T}) \in \mathbf{Z}[\mathrm{T}_1,..,\mathrm{T}_k]\) and integers \(s_1 \geqslant 1, ..., s_k \geqslant 1\) such that we have

\[
\sum_ {n \in \mathbb {N} ^ {k}} \operatorname{long} _ {\mathrm{A}} (\mathrm{H} _ {n}). \mathrm{T} ^ {n} = \frac {\mathrm{Q} (\mathrm{T})}{(1 - \mathrm{T} _ {1}) ^ {s _ {1}} \dots (1 - \mathrm{T} _ {k}) ^ {s _ {k}}}.
\]

(We proceed by induction on d.)

\begin{enumerate}
\def\labelenumi{\alph{enumi})}
\setcounter{enumi}{2}
\tightlist
\item
  From this, deduce, by extending to \(\mathbf{Z}((\mathrm{T}_{1},\ldots,\mathrm{T}_{k}))\) the results of § 6, no. 1, the existence of integers \(c(\alpha_{1},\ldots,\alpha_{k})\) for the sequences \((\alpha_{1},\ldots,\alpha_{k})\) such that \(|\alpha| = \sum_{i=1}^{k}\alpha_{i} = d\) such that we have
\end{enumerate}

\[
e _ {\mathfrak {q} _ {1} ^ {\nu_ {1}} \dots \mathfrak {q} _ {k} ^ {\nu_ {k}}} (\mathbf {M}) = \sum_ {| \alpha | = d} \frac {d !}{\alpha_ {1} ! \dots \alpha_ {k} !} c (\alpha_ {1},..., \alpha_ {k}) v _ {1} ^ {\alpha_ {1}}... v _ {k} ^ {\alpha_ {k}}
\]

for \((\mathbf{v}_1,\dots,\mathbf{v}_k)\) in \(\mathbf{N}^k\)

\begin{enumerate}
\def\labelenumi{\alph{enumi})}
\setcounter{enumi}{3}
\tightlist
\item
  Suppose now that A is local, with infinite residue field κ. Show, using the existence of superficial elements for the graded H-module N, that one has
\end{enumerate}

\[
c (\alpha_ {1}, \dots , \alpha_ {k}) = \inf _ {\mathbf {q} _ {\alpha}} (e _ {\mathbf {q} _ {\alpha}} (\mathbf {M}))
\]

where \(q_{\alpha}\) ranges over the set of ideals of A generated by \(\alpha_{1}\) elements of \(q_{1}\), \(\alpha_{2}\) elements of \(q_{2}\), \ldots, \(\alpha_{k}\) elements of \(q_{k}\). We sometimes denote \(e(\mathbf{q}_{1}^{[\alpha_{1}]} + \cdots + \mathbf{q}_{k}^{[\alpha_{k}]}; \mathbf{M})\) by the integer \(c(\alpha_{1}, \ldots, \alpha_{k})\), abbreviated as \(e(\mathbf{q}_{1}^{[\alpha_{1}]} + \cdots + \mathbf{q}_{k}^{[\alpha_{k}]})\) when M = A. e) Prove the equality

\[
e _ {\mathfrak {q} _ {1}, \mathfrak {q} _ {2}} (\mathbf {M}) = \sum_ {i = 1} ^ {d} \binom {d} {i} e (\mathfrak {q} _ {1} ^ {[ i ]} + \mathfrak {q} _ {2} ^ {[ d - i ]}; \mathbf {M}).
\]

\begin{enumerate}
\def\labelenumi{\alph{enumi})}
\setcounter{enumi}{5}
\tightlist
\item
  Show that we have:
\end{enumerate}

\[
e \left(\overline {{{\mathfrak {q}}}} _ {1} ^ {[ \alpha_ {1} ]} + \dots + \overline {{{\mathfrak {q}}}} _ {k} ^ {[ \alpha_ {k} ]}; \mathrm{M}\right) = e \left(\mathfrak {q} _ {1} ^ {[ \alpha_ {1} ]} + \dots + \mathfrak {q} _ {k} ^ {[ \alpha_ {k} ]}; \mathrm{M}\right)
\]

where \(\overline{\mathfrak{q}}_{i}\) denotes the integral closure in A of the ideal \(\mathfrak{q}_{i}\) (cf. exercise 19).

\begin{enumerate}
\def\labelenumi{\arabic{enumi})}
\setcounter{enumi}{20}
\tightlist
\item
  Let A and A' be two Noetherian rings, and let q (resp. q') be an ideal of A (resp. A') contained in the radical of A (resp. A') and such that \(\operatorname{long}(A/q)\) (resp. \(\operatorname{long}(A'/q'))\) is finite. Let Q be the ideal \(q \otimes A' + A \otimes q'\) of the ring \(A \otimes_{\mathbf{Z}} A'\). Show that it is contained in the radical of \(A \otimes_{\mathbf{Z}} A'\) and that one has the equality
\end{enumerate}

\[
e _ {\mathfrak {Q}} (\mathrm{A} \otimes_ {\mathbf {Z}} \mathrm{A} ^ {\prime}) = e _ {\mathfrak {q}} (\mathrm{A}). e _ {\mathfrak {q} ^ {\prime}} (\mathrm{A} ^ {\prime}).
\]

22) * Let A be a Noetherian integral domain, which is local and complete (resp. a finitely generated algebra over a field k), \(a \in A\) a non-invertible and nonzero element. We make the following assumptions:

α) \(\operatorname{Supp}(A/aA)\) has a single minimal element p.

\(\beta)\) We have \(aA_{p} = pA_{p}\).

γ) A/p is an integrally closed ring.

The aim of this exercise is to show that, for every maximal ideal m of A containing a, the ring \(A_{m}\) is integrally closed and that one has p = aA ("Hironaka's lemma"). In this exercise we shall use the catenarity property. This property is satisfied by integral, Noetherian, local, complete rings, as will be seen in Chapter X. It is also satisfied by integral algebras of finite type over a field (§ 2, n° 4, th. 3).

\begin{enumerate}
\def\labelenumi{\alph{enumi})}
\item
  Show that \(A_p\) is a discrete valuation ring (§ 3, n° 1, prop. 1 and VI, § 3, n° 6, prop. 9).
\item
  We make the hypotheses \(\alpha\) and \(\beta\) and suppose A is integrally closed. Show that one has \(\mathfrak{p} = a\mathbf{A}\) (VII, § 1, n° 4, prop. 8).
\item
  Let A' be the integral closure of A. Prove that A' is an integral, Noetherian, local and complete ring (resp. that it is an algebra of finite type over k). (Note that A is Japanese (IX, § 4, n° 2, th. 2), hence A' is integral, Noetherian, semi-local and complete and, consequently, local (III, § 2, n° 13, corollary to prop. 19). In the case where A is an algebra of finite type over k, use V, § 3, n° 2, th. 2.)
\item
  Show that one has A' ⊗A A\_p = A\_p (use a) and V, § 1, n° 5, prop. 16). Deduce that there exists a single prime ideal p' of A' lying over p and that one has A\_p = A'\_p', aA'\_p' = p'A'\_p'.
\item
  Let q be a minimal element of the support of A'/aA'. Show that q is a prime ideal of A' lying over p. (Note that q has height 1 (§ 3, n° 1, prop. 1) and that consequently, by virtue of catenarity, one has \(\dim(A'/q)=\dim(A)-1\) . Deduce that one has \(\dim(A/(q \cap A))=\dim(A)-1\) and consequently that \(q \cap A = p\) .)
\item
  Deduce from e) that \(aA'\) is a prime ideal of \(A'\) and that one has \(aA' = p' = pA'\) (one may use b)).
\end{enumerate}

g) Show that one has \(\mathbf{A}_{\mathfrak{m}}^{\prime} = \mathbf{A}_{\mathfrak{m}}\) for every maximal ideal m of A containing \(a\). (From the inclusions \(\mathbf{A} \subset \mathbf{A}' \subset \mathbf{A}_{\mathfrak{p}}\), deduce the inclusions \(\mathbf{A}/\mathfrak{p} \subset \mathbf{A}'/\mathfrak{p}' \subset \mathbf{A}_{\mathfrak{p}}/\mathfrak{p}\mathbf{A}_{\mathfrak{p}}\). Deduce that \(\mathbf{A}/\mathfrak{p}\) and \(\mathbf{A}'/\mathfrak{p}'\) have the same field of fractions and consequently, according to \(\gamma)\), that \(\mathbf{A}/\mathfrak{p} = \mathbf{A}'/\mathfrak{p}' = \mathbf{A}' \otimes_{\mathbf{A}} (\mathbf{A}/\mathfrak{p})\), the last equality resulting from \(f)\). One therefore has \((\mathbf{A}'/\mathbf{A}) \otimes_{\mathbf{A}} (\mathbf{A}/\mathfrak{p}) = 0\). Conclude.) Then complete the proof of Hironaka's lemma. *

\begin{enumerate}
\def\labelenumi{\arabic{enumi})}
\setcounter{enumi}{22}
\tightlist
\item
  \begin{enumerate}
  \def\labelenumii{\alph{enumii})}
  \tightlist
  \item
    Let A be a Noetherian local ring such that \(e(A) = 1\). Show that A possesses a single minimal prime ideal \(p\) such that \(\dim(A) = \dim(A/p)\). Show moreover that one has \(\text{long}(A_p) = 1\) and \(e(A/p) = 1\) (\S 7, no. 1, remark 4).
  \end{enumerate}
\end{enumerate}

\begin{enumerate}
\def\labelenumi{\alph{enumi})}
\setcounter{enumi}{1}
\tightlist
\item
  Recall (p. 82, exerc. 10) that a local ring is equidimensional if, for every minimal prime ideal q of A, one has \(\dim(A/q) = \dim(A)\), and that it has no embedded prime ideals if the A-module A has no embedded associated prime ideals (IV, § 2, no. 3, remark). * (It will be seen in Chapter X that the completions of integral quotients of Macaulay local rings possess these two properties.) * Let A be a local, Noetherian, equidimensional ring with no embedded prime ideals, such that \(e(A) = 1\). Show that A is integral.
\end{enumerate}

24) * Let A be a Noetherian local ring, \(\widehat{A}\) its completion. The aim of this exercise is to prove the equivalence of the properties:

\begin{enumerate}
\def\labelenumi{(\roman{enumi})}
\item
  A is regular;
\item
  \(\hat{\mathbf{A}}\) is equidimensional and has no embedded prime ideals, and \(e(\mathbf{A}) = 1\) ;
\item
  \(\hat{\mathbf{A}}\) is integral and \(e(\mathbf{A}) = 1\) .
\end{enumerate}

Prove the implications (i) \(\Rightarrow\) (ii) \(\Rightarrow\) (iii). Note that (iii) implies the equality \(e(\hat{\mathbf{A}}) = 1\) and that, to prove (iii) \(\Rightarrow\) (i), it suffices to prove that \(\hat{\mathbf{A}}\) is regular. One may therefore, to prove (iii) \(\Rightarrow\) (i), suppose A complete, which we shall do from now on. We shall first treat the case where the residue field \(\kappa_{\mathrm{A}}\) has infinitely many elements. One proceeds by induction on the dimension of A. Examine the case \(\dim(A) = 0\), and henceforth suppose \(\dim(A) > 0\). There then exists a superficial element \(x \in \mathrm{A}\) (\S 7, no. 5, remark 4). By virtue of the catenarity property (exerc. 22) and of \S 3, no. 1, prop. 1, for every minimal prime ideal p among those containing xA, one has \(\dim(A/p) = \dim(A/xA)\). Since one has \(e(\mathrm{A}/x\mathrm{A}) = 1\) (\S 7, no. 5, th. 1), there exists a single minimal prime ideal p among those containing xA, and one has \(xA_{p} = pA_{p}\), \(e(\mathrm{A}/p) = 1\) (exerc. 23). By the induction hypothesis, A/p is regular, hence integrally closed (\S 5, no. 2, cor. 1 to th. 1). By Hironaka's lemma (exerc. 22), one deduces p = xA. Hence A/xA is regular and since A is integral, A is regular (\S 5, no. 3, prop. 2).

Suppose now that \(\kappa_{A}\) is a finite field. Prove that it suffices to show that the blowing-up A{]}X{[} (IX, Appendix, no. 2) is regular. Note that A and consequently A{]}X{[} are isomorphic to quotients of a regular ring (IX, § 2, no. 5, th. 3), and \(e(A]X[)=1\). By the theory of Macaulay rings (Chapter X), the completion \(\widehat{A]X[}\) is equidimensional and has no embedded prime ideals. Deduce that \(\widehat{A]X[}\) is integral (exerc. 23) and that \(e(\widehat{A]X[})=1\), hence that \(A]X[\) is regular. Conclude. *

\begin{enumerate}
\def\labelenumi{\arabic{enumi})}
\setcounter{enumi}{24}
\tightlist
\item
  Let \(k\) be a field, A a \(k\) -local algebra, localized at a prime ideal of a \(k\) -algebra of finite type. Show that A is regular if and only if A is integral and \(e(A) = 1\) . (One may be inspired by the method followed in exerc. 24 or also note that since A is an integral quotient of a regular ring, the local ring \(\hat{A}\) is equidimensional and has no embedded prime ideals.)
\end{enumerate}

